\section{Introduction}

The question of whether large language models (LLMs) possess ``personality'' has moved from philosophical speculation to empirical investigation. A growing body of work has administered human personality questionnaires---the Big Five Inventory, the Dark Triad scales, and others---to LLMs, reporting personality profiles that appear coherent on the surface \citep{miotto2023gpt3, bhandari2025evaluating, serapio2025psychometric}. These studies typically present domain-level scores, rank models along trait dimensions, and conclude that LLMs exhibit measurable personality differences.

However, this research program rests on an untested assumption: that human-validated psychometric instruments retain their measurement properties when applied to non-human respondents. In personality psychology, an instrument is only valid for a population if it demonstrates adequate internal consistency (Cronbach's $\alpha$), recovers the expected factor structure, shows convergent validity with theoretically related constructs, and passes measurement invariance tests across subgroups \citep{cronbach1951coefficient, lord1968statistical}. None of these checks are routinely performed before interpreting LLM personality scores.

Recent work has begun to surface cracks. \citet{suhr2025challenging} showed that LLMs frequently endorse both forward-coded and reverse-coded items simultaneously, violating basic internal consistency. \citet{acerbi2024llm} demonstrated that LLMs shift their responses toward socially desirable endpoints by up to 1.20 human standard deviations. \citet{pelet2025persistent} found that LLM personality measurements are persistently unstable across question reorderings, model scales, and conversation histories. Together, these findings suggest that LLM personality scores may reflect response artifacts rather than genuine trait variation.

We take a different approach: rather than profiling LLM personality, we ask whether the instruments themselves are valid for this population. We administer a 221-item psychometric battery---comprising the IPIP-NEO-120 \citep{johnson2014measuring}, the Short Dark Triad \citep{jones2014introducing}, the ZKPQ-50-CC \citep{aluja2006cross}, and the EPQR-A \citep{francis1992development}---to 18 LLMs across 17 role-conditioned prompts, and subject the resulting 67,626 responses to six standard psychometric validation analyses. Our central finding is that human personality instruments do not transport cleanly to LLMs: internal consistency collapses ($\alpha$ ranges from $-0.02$ to $0.36$, versus $0.87$--$0.90$ in human norms), the Big Five factor structure reduces from five to three factors, and inter-model variance accounts for less than 1\% of response variability.


\section{Related Work}

\subsection{Personality Measurement in LLMs}

The application of personality instruments to LLMs has grown rapidly. \citet{miotto2023gpt3} administered the Big Five Inventory (BFI) to GPT-3, reporting measurable personality profiles. \citet{serapio2025psychometric} developed a comprehensive psychometric framework for evaluating and shaping personality traits across 18 LLMs, finding that instruction-tuned models can produce seemingly reliable measurements under specific prompt conditions. \citet{bhandari2025evaluating} applied five personality measurement tools to multiple LLMs, reporting varying consistency across personality dimensions. However, these studies predominantly report descriptive personality profiles without validating that the instruments measure coherent constructs in LLMs. The critical step---establishing measurement invariance between human and LLM populations---is typically assumed rather than tested \citep{suhr2025challenging}.

\subsection{Response Style and Acquiescence Bias}

Survey methodology has long recognized that self-report data is contaminated by response styles---systematic tendencies to respond to items regardless of content. Acquiescence bias (the tendency to agree with statements regardless of direction) is particularly damaging for instruments that use reverse-coded items as an internal consistency check \citep{lord1968statistical}. \citet{acerbi2024llm} demonstrated that LLMs exhibit human-like social desirability biases in Big Five surveys, shifting scores toward socially desirable endpoints. Critically, they found this bias persists even after item randomization and rephrasing, and they ruled out pure acquiescence as the sole mechanism. This suggests that LLM response contamination is multi-determined, combining acquiescence with social desirability and potentially alignment-driven prosociality.

\subsection{Psychometric Validation Methods}

Classical test theory provides a well-established toolkit for validating measurement instruments. Cronbach's $\alpha$ \citep{cronbach1951coefficient} quantifies internal consistency within a scale. Exploratory factor analysis with the Kaiser criterion \citep{kaiser1960application} tests whether the expected latent structure emerges from the data. Tucker's congruence coefficient \citep{tucker1951method} quantifies factor similarity across populations. Measurement invariance testing establishes whether an instrument measures the same construct across groups---a prerequisite for any meaningful cross-population comparison. These methods have been applied to the IPIP-NEO-120 \citep{johnson2014measuring}, the SD3 \citep{jones2014introducing}, and the ZKPQ-50-CC \citep{zuckerman2002zuckerman} in human populations, yielding well-established benchmarks: Cronbach's $\alpha$ typically ranges from $0.87$ to $0.90$ for Big Five domains, and the five-factor structure is consistently recovered with Tucker's $\phi > 0.95$ \citep{mccrae2005universal}. To our knowledge, no study has systematically applied this full validation battery to LLM personality data.

\subsection{Alignment Effects on LLM Behavior}

RLHF and safety alignment training fundamentally shape LLM behavior. \citet{ouyang2022training} showed that reinforcement learning from human feedback produces models that are more helpful and harmless, but potentially less diverse in their outputs. \citet{bai2022constitutional} introduced Constitutional AI as an approach to training safer models through principle-guided self-correction. These alignment procedures may create a universal ``moderate prosocial'' response profile that overrides genuine variation across models and training regimes. If alignment homogenizes responses, then personality differences between models become artifacts of residual noise rather than meaningful trait variation.
